%
\hsection{Classes / Klassen}%
%
%
\begin{question}{What is it? / Was ist das?}%
\begin{itemize}%
\item[EN] Show the syntax of classes in \python. Include and additionally explain the following elements\begin{enumerate}%
\item the class name,
\item the initializer method,
\item the attributes as well as how and where they declared/initialized,
\item methods of a class,
\item the use of \pythonil{self}.
\end{enumerate}%
%
\item[DE] Zeigen Sie die Syntax von Klassen in \python. Beinhalten und erklären Sie die folgenden Elemente\begin{enumerate}%
\item den Name der Klasse,
\item die Initialisierer-Methode,
\item die Attribute sowohl wie und wo sie deklariert/initialisiert werden,
\item Methoden,
\item die Verwendung von \pythonil{self}.
\end{enumerate}
\end{itemize}%
\end{question}%
%
%
\begin{question}{Use Case? / Nutzen?}%
\begin{itemize}%
\item[EN] Name and explain two general use cases of classes.%
%
\item[DE] Nennen und erklären Sie zwei generelle Szenarien, in denen man Klassen verwenden kann/sollte.%
\end{itemize}%
\end{question}%
%
%
\begin{question}{What is it? / Was ist das?}%
\begin{itemize}%
\item[EN] Explain what \emph{immutable} class (instances) are:\begin{enumerate}%
\item When is a class instance immutable?
\item What do we need to do in \python\ so that a class instance is annotated to be immutable?%
\item Give two advantages of making classes immutable.%
\end{enumerate}%
%
\item[DE] Erklären Sie, was \emph{unveränderliche} Klassen(-Instanzen) sind:\begin{enumerate}
\item Wann ist eine Klassen(-Instanz) unveränderlich?
\item Was What do we need to do in \python\ so dass eine Klasse-Instanz als unveränderlich annotiert ist?%
\item Nennen Sie zwei Vorteile von unveränderlichen Klassen.%
\end{enumerate}
\end{itemize}%
\end{question}%
%
%
%
\begin{question}{Syntax}%
\begin{itemize}%
\item[EN] In which situation would we start the name of an attribute of a class with two underscore characters, e.g., \pythonil{__counter}?%
%
\item[DE] In welcher Situation würden wir den Name eines Attributs einer Klasse mit zwei Unterstrichen beginnen, zum Beispiel~\pythonil{__counter}?%
\end{itemize}%
\end{question}%
%
%
%
%
\begin{question}{Programming / Programmieren}%
\gitLoadPython{class_01}{}{questions/class_01.py}{}%
\gitExecPython{class_01}{}{questions}{class_01.py}{}%
\listingPythonAndOutput{class_01}{The program \programUrl{class_01}}{}%
\begin{itemize}%
\item[EN] The output of program \programUrl{class_01} is given in \cref{exec:class_01} above. %
It imports a class \pythonil{Person} from another module.
Write a class \pythonil{Person} in \python\ that has the necessary attributes, initializer, and method to produce exactly the observed behavior and output.%
\item[DE] Die Ausgabe des Programms \programUrl{class_01} ist oben in \cref{exec:class_01} gegeben. %
Es importiert die Klasse \pythonil{Person} aus einem anderen Modul.
Schreiben Sie die Klasse \pythonil{Person} in \python\ mit den notwendigen Attributen, den notwendigen Initialisierer und der notwendigen Methode so dass genau das beobachtete Verhalten bzw.\ die beobachtete Ausgabe entsteht.%
\end{itemize}%
\end{question}%

%
\begin{question}{Programming / Programmieren}%
\begin{align}%
z_1+z_2=(a_1+b_1\imaginaryUnit)+(a_2+b_2\imaginaryUnit)&=(a_1+a_2)+(b_1+b_2)\imaginaryUnit\\%
z_1-z_2=(a_1+b_1\imaginaryUnit)-(a_2+b_2\imaginaryUnit)&=(a_1-a_2)+(b_1-b_2)\imaginaryUnit\\%
z_1*z_2=(a_1+b_1\imaginaryUnit)*(a_2+b_2\imaginaryUnit)&=(a_1a_2 - b_1b_2) + (a_1b_2 + a_2b_1)\imaginaryUnit\\%
\complexConj{z_1}=\complexConj{a_1+b_1\imaginaryUnit}&=a_1-b_1\imaginaryUnit\quad\text{(\emph{conjugate})}%
\end{align}%
\begin{itemize}%
\item[EN] The above equations demonstrate mathematical operations for one to two complex numbers~$z_1,z_2\in\complexNumbers$ such that $z_1=a_1+b_1\imaginaryUnit$ and $z_2=a_2+b_2\imaginaryUnit$, where \imaginaryUnit\ is the imaginary unit with~$\imaginaryUnit^2=-1$~\cite{VDC2022AITCN}.\begin{enumerate}%
\item Write a \python\ class for representing complex numbers.%
\item The instances of the class should be immutable.%
\item The class should implement the four operations \pythonil{plus}, \pythonil{minus}, \pythonil{multiply}, and \pythonil{conjugate} as defined by the above equations.%
\item The class must be annotated with a \pgls{docstring} which contains a \pythonil{doctest}.%
\item All attributes of the class must be annotated with \pglspl{typeHint}.%
\end{enumerate}%
%
\item[DE] Die Gleichungen oben stellen Operationen für ein bis zwei komplexe Zahlen~$z_1,z_2\in\complexNumbers$ mit $z_1=a_1+b_1\imaginaryUnit$ und $z_2=a_2+b_2\imaginaryUnit$ dar, wobei \imaginaryUnit\ die imaginäre Einheit mit~$\imaginaryUnit^2=-1$ ist~\cite{VDC2022AITCN}.\begin{enumerate}%
\item Schreiben Sie eine \python-Klasse um komplexe Zahlen darstellen zu können.%
\item Die Instanzen der Klasse sollen unveränderlich sein.%
\item Die Klasse soll die vier Operationen \pythonil{plus}, \pythonil{minus}, \pythonil{multiply} und \pythonil{conjugate} wie in den Gleichungen oben definiert implementieren.%
\item Die Klasse soll mit einem DocString annotiert werden, der einen DocTest beinhaltet.%
\item Alle Attribute der Klasse müssen mit Type-Hints annotiert werden.%
\end{enumerate}%
\end{itemize}%
\end{question}%
%
%
\begin{question}{Programming / Programmieren}%
\begin{align}%
q_1+q_2&=\frac{a_1}{b_1}+\frac{a_2}{b_2}=\frac{a_1b_2+a_2b_1}{b_1b_2}~\quad\text{reduce/kürzen!}\\%
q_1-q_2&=\frac{a_1}{b_1}-\frac{a_2}{b_2}=\frac{a_1b_2-a_2b_1}{b_1b_2}~\quad\text{reduce/kürzen!}\\%
q_1*q_2&=\frac{a_1}{b_1}*\frac{a_2}{b_2}=\frac{a_1a_2}{b_1b_2}~\quad\text{reduce/kürzen!}\\%
q_1/q_2&=\frac{a_1}{b_1}/\frac{a_2}{b_2}=\left\{\begin{array}{rl}%
\pythonil{ZeroDivisionError}&\text{~if/wenn $q_2=0$}\\%
\frac{a_1b_2}{b_1a_2}&\text{~otherwise/else}%
\end{array}\right.~\quad\text{reduce/kürzen!}%
\end{align}%
\begin{itemize}%
\item[EN] The above equations demonstrate mathematical operations for two rational numbers~$q_1,q_2\in\rationalNumbers$ such that $q_1=\frac{a_1}{b_1}$ and $q_2=\frac{a_2}{b_2}$.
In other words, $q_1$ and $q_2$ are fractional numbers where $a_1,a_2,b_1,b_2\in\integerNumbers$~(integer numbers) and $b1,b2\neq0$.\begin{enumerate}%
\item Write a \python\ class for representing rational numbers.%
\item The class should properly represent~0 and negative numbers.%
\item The class should always ensure that the fractions are always reduced as far as possible.
In other words, with the exception of~0, it holds that $gcd(a,b)=1$ for all~$\frac{a}{b}$.
(You can use the function~\pythonil{gcd} from the module~\pythonil{math}.)%
\item The instances of the class should be immutable.%
\item The class should implement the four operations \pythonil{plus}, \pythonil{minus}, \pythonil{multiply}, and \pythonil{divide} as defined by the above equations.%
\item A division by~0 should raise an \pythonil{ZeroDivisionError}.
\item The class must be annotated with a \pgls{docstring} which contains a \pythonil{doctest}.%
\item All attributes of the class must be annotated with \pglspl{typeHint}.%
\end{enumerate}%
%
\item[DE] Die Gleichungen oben demonstrieren die mathematischen Operationen für zwei rationale Zahlen~$q_1,q_2\in\rationalNumbers$ mit $q_1=\frac{a_1}{b_1}$ und $q_2=\frac{a_2}{b_2}$.
Mit anderen Worten, $q_1$ und $q_2$ sind Brüche mit $a_1,a_2,b_1,b_2\in\integerNumbers$~(sind Ganzzahlen) und $b_1,b_2\neq0$.\begin{enumerate}%
\item Schreiben Sie eine \python-Klasse, um rationale Zahlen darzustellen.%
\item Die Klasse sollte sowohl 0 als auch negative Zahlen korrekt darstellen.%
\item Die Klasse sollte die Brüche immer so weit wie möglich gekürzt darstellen. %
Also mit der Ausnahme von 0 gilt immer das $gcd(a,b)=1$ für alle~$\frac{a}{b}$. %
(Sie können die Funktion~\pythonil{gcd} aus dem Modul~\pythonil{math} verwenden.)%
\item Die Instanzen der Klasse sollen unveränderlich sein.%
\item Die Klasse soll die vier Operationen \pythonil{plus}, \pythonil{minus}, \pythonil{multiply} und \pythonil{divide} wie in den Gleichungen oben definiert implementieren.%
\item Eine Division durch~0 sollte einen \pythonil{ZeroDivisionError} auslösen.
\item Die Klasse soll mit einem DocString annotiert werden, der einen DocTest beinhaltet.%
\item Alle Attribute der Klasse müssen mit Type-Hints annotiert werden.%
\end{enumerate}%
\end{itemize}%
\end{question}%
%
%
%
\begin{question}{Programming / Programmieren}%
\begin{itemize}%
\item[EN] Let's build a class hierarchy for life forms. Write code for everything below.%
\begin{enumerate}%
\item The base class be \pythonil{LifeForm}, and it has one attribute~\emph{name}, which represents the name of a life form as string.%
\item The direct subclasses of \pythonil{LifeForm} are \pythonil{Plant} and \pythonil{Animal}.%
\item The class \pythonil{Plant} has a method \pythonil{grow}, which prints \pythonil{\"(name) grows!\"}, where \textil{(name)} is to be replaced by the value of the \emph{name} attribute.%
\item The class \pythonil{Animal} has a method \pythonil{eat}, which prints \pythonil{\"(name) eats!\"}, where \textil{(name)} is to be replaced by the value of the \emph{name} attribute.%
\item Write code that instantiates \pythonil{Plant} and creates a plant named \textil{Willi}.
Call its \pythonil{grow} method.%
\item Create a subclass of \pythonil{Animal} named \pythonil{Lion}.
This subclass gets a method \pythonil{do_something(hungry: bool)}, which invokes the \pythonil{eat} method if \pythonil{hungry == True}.%
\end{enumerate}%
%
\item[DE] Lassen Sie uns eine Klassenhierarchie für Lebensformen bauen. Schreiben Sie Kode für \emph{alles} Folgendes:%
\begin{enumerate}%
\item Die Basisklasse ist \pythonil{LifeForm} und sie hat ein Attribut \emph{name}, welches den Name der Lebensform als String repräsentiert.%
\item Die direkten Subklassen von \pythonil{LifeForm} sind \pythonil{Plant}~\inDE{Pflanze} und \pythonil{Animal}~\inDE{Tier}.%
\item Die Klasse \pythonil{Plant} hat eine Methode \pythonil{grow}, die ausgibt \pythonil{\"(name) grows!\"}, wobei \textil{(name)} mit dem Wert der \emph{name}-Attributs zu ersetzen ist.%
\item Die Klasse \pythonil{Animal} hat eine Methode \pythonil{eat}, die ausgibt \pythonil{\"(name) eats!\"}, wobei \textil{(name)} mit dem Wert des Attributs \emph{name} zu ersetzen ist.%
\item Schreiben Sie Kode der die Klasse \pythonil{Plant} instantiiert und eine Pflanze namens \textil{Willi} erzeugt.
Rufen Sie deren \pythonil{grow}-Methode auf.%
\item Erstellen Sie eine Subklasse von \pythonil{Animal}, die Sie \pythonil{Lion}~\inDE{Löwe} nennen.
Diese Subklasse bekommt eine Methode \pythonil{do_something(hungry: bool)}, die die \pythonil{eat}-Methode aufruft wenn \pythonil{hungry == True}.%
\end{enumerate}%
\end{itemize}%
\end{question}%
%
\endhsection%
%
